% Table T24 -- Complexity Analysis
% Experiment: EXP-A2
% Objective: O5 (complexity and deployability)
% Source: D:/Projects/HeartGuard/outputs/11_complexity/training_time_summary.csv; D:/Projects/HeartGuard/outputs/11_complexity/inference_stage_summary.csv; D:/Projects/HeartGuard/outputs/11_complexity/model_footprint.csv; D:/Projects/HeartGuard/outputs/11_complexity/complexity_overview.csv
% Generated by PV-MEPCG / PulseVision at 2026-09-10T11:37:47.956231+00:00
\begin{table}[htbp]
  \centering
  \caption{Computational cost beside the performance it buys, one row per model. 'Bench fit' and every size and memory column come from one controlled fit on binary fold r0f0 under configs/models.yaml defaults, on an otherwise idle machine. 'Run fit mean' and 'Run fit max/min' come from the EXP-A2 run's own instrumentation across its folds and are therefore a different quantity: a searched configuration, measured while the project was running. Metrics are that run's fold means.}
  \label{tab:t24}
  \begin{tabular}{lrrrrrrrrrrrrr}
    \toprule
    Model & Task & Benched configuration & Features & Bench fit (s) & Run fit mean (s) & Run fit max/min & Predict, 1 record (s) & Predict, batched (s/record) & On disk (MB) & Peak RSS above baseline (MB) & Whole pipeline (s/recording) & balanced accuracy & sensitivity \\
    \midrule
    M1 & binary & configs/models.yaml defaults & 138 & 0.49 & 2.97 & 320.0 & 0.001405 & 0.000007 & 0.007 & 12.9 & 2.343 & 0.847 & 0.865 \\
    M2 & binary & configs/models.yaml defaults & 138 & 0.09 & n/a & n/a & 0.006245 & 0.003180 & 2.889 & 14.3 & 2.343 & n/a & n/a \\
    M3 & binary & configs/models.yaml defaults & 138 & 1.11 & 1.40 & 2.3 & 0.007671 & 0.000487 & 2.214 & 15.6 & 2.343 & 0.800 & 0.667 \\
    M4 & binary & configs/models.yaml defaults & 138 & 14.50 & 13.00 & 23.2 & 0.177161 & 0.000349 & 12.471 & 24.6 & 2.343 & 0.852 & 0.819 \\
    M5 & binary & configs/models.yaml defaults & 138 & 28.77 & 99.18 & 15.9 & 0.001408 & 0.000013 & 0.273 & 15.1 & 2.343 & 0.847 & 0.813 \\
    M6 & binary & configs/models.yaml defaults & 138 & 228.46 & 545.14 & 9.4 & 0.182914 & 0.000896 & 15.562 & 111.5 & 2.343 & 0.856 & 0.862 \\
    M7 & binary & configs/models.yaml defaults & 138 & 257.23 & 534.58 & 10.5 & 0.169259 & 0.000859 & 15.562 & 120.3 & 2.343 & 0.856 & 0.862 \\
    M8 & binary & configs/models.yaml defaults & 138 & 6.63 & 12.94 & 42.3 & 0.003835 & 0.000019 & 0.673 & 51.2 & 2.343 & 0.848 & 0.822 \\
    \bottomrule
  \end{tabular}
  \par\footnotesize THE MODEL IS NOT THE LATENCY. End-to-end inference on a real recording is dominated by feature extraction, which is identical for every model; the 'Whole pipeline' column is the same value in every row for that reason. A 190x difference in predict time is invisible to anyone waiting for a result.
  \par\footnotesize 'Peak RSS above baseline' is the process resident set sampled every 20 ms during the fit, minus the reading taken immediately before it. It is a peak DURING the fit, not a peak CAUSED by it, and a fit shorter than a few sample intervals carries too few samples to mean anything -- model\_footprint.csv reports the sample count and a memory\_measurement\_reliable flag per row.
  \par\footnotesize The two fit-time columns are not comparable with each other and neither is an error. See T25.
  \par\footnotesize Measured on CPU only -- this machine has no CUDA GPU, and no number here may be compared against a GPU-trained model. M9 (1D-CNN) is out of scope for that reason and is not fitted anywhere in this project.
  \par\footnotesize PV-MEPCG / PulseVision is an academic screening and decision-support prototype, not a diagnostic tool.
\end{table}
